% =============================================================================
%  5. Results
%  The subsection ORDER is the argument: each one removes an objection to the
%  one before it. 5.1 is the claim; 5.2 answers "you picked a weak judge";
%  5.3 answers "is any of this needed?"; 5.4 answers "private data only";
%  5.5 answers "why not one model per camera?".
% =============================================================================
\section{Results}
\label{sec:results}

\subsection{Four proposal stages under a frozen judge}
\label{sec:results:main}

\tabref{tab:main} sweeps the \proposer{} with the \judge{} held fixed.

\begin{table}[t]
\centering
\caption{Four proposal stages, judge held fixed. 68 cases, 73 instances, three
trams, nine views. Object recall does not move; everything else does.}
\label{tab:main}
\small
\setlength{\tabcolsep}{4.5pt}
\begin{tabular}{@{}lrrrrrrr@{}}
\toprule
\textbf{Proposal stage} & \textbf{Regions} & \textbf{Frame $F_1$} &
\textbf{Frame Re} & \textbf{Frame Sp} & \textbf{Obj.\ recall} &
\textbf{Strict} & \textbf{Reg.\ prec.} \\
\midrule
\armphoto{}          & 1192 & 0.931 & 0.871 & 1.000 & 55/73 & 0.534 & 0.694 \\
\armphotodino{}      &  727 & 0.931 & 0.871 & 1.000 & 55/73 & 0.534 & 0.773 \\
\armdino{}           & 1228 & \textbf{0.984} & \textbf{0.968} & 1.000 & 55/73 & \textbf{0.630} & \textbf{0.896} \\
\armdinodinomaly{}   & \textbf{654} & \textbf{0.984} & \textbf{0.968} & 1.000 & 55/73 & \textbf{0.630} & \textbf{0.896} \\
\bottomrule
\end{tabular}
\end{table}

The object-recall column is the one to read first: it does not move.
Fifty-five of seventy-three instances, in all four arms, to the instance.
Everything else does move --- frame recall from 0.871 to 0.968, region
precision from 0.694 to 0.896, kept-but-wrong regions from 30 down to 7, and
strict-IoU boxing from 0.534 to 0.630 --- while frame specificity stays at
1.000 throughout, so none of this is bought with false alarms.

Three consequences follow, and the third is the paper.

First, \emph{``better boxes find more objects'' is the tempting reading of
\tabref{tab:main} and it is false.} The localization-only view
(\tabref{tab:localization}) shows the feature \proposer{} localizing four
instances the photometric one misses --- and the \judge{} rejects them anyway.
The \proposer{} sets the ceiling; the \judge{} decides where under it the
system sits.

Second, \emph{a veto is not a \proposer{}.} Rows 1--2 of \tabref{tab:main} are
identical on every quality column and differ only in cost, because a veto
deletes boxes and never redraws one; the same holds for rows 3--4, where
\armdinodinomaly{} matches \armdino{} to the digit for 47\,\% fewer judge
calls. Matching the ceiling means none of the 574 deleted boxes held a true
instance. So the operative question is never \emph{which filter} but \emph{who
proposes}.

Third, \emph{the binding constraint is spatial, not linguistic.} With the
\judge{} frozen, changing the \proposer{} moves frame $F_1$ by 0.053 and region
precision by 0.202. \secref{sec:results:judge} shows what changing the
\judge{} does instead.

\begin{table}[t]
\centering
\caption{Localization only, no judge calls. The lenient rule barely separates
the two families; the strict rule separates them decisively.}
\label{tab:localization}
\begin{tabular}{@{}lrrr@{}}
\toprule
\textbf{Proposal stage} & \textbf{Instances localized} & \textbf{Strict IoU} & \textbf{Regions} \\
\midrule
\armphoto{}             & 65/73 & 42/73 & 1192 \\
\armphotodino{}         & 65/73 & 42/73 &  727 \\
Dinomaly alone          & 62/73 & 48/73 &  813 \\
\armdino{}              & \textbf{69/73} & \textbf{58/73} & 1228 \\
\armdinodinomaly{}      & \textbf{69/73} & \textbf{58/73} &  654 \\
\bottomrule
\end{tabular}
\end{table}

The mechanism behind the strict/lenient gap is blunt and measurable. On tram
3333 the photometric \proposer{} scores 5/26 strict against the feature
\proposer{}'s 19/26, and its worst box covers $911\,360$~px of a
$1280 \times 720$ frame --- 98.9\,\% of the image. Such a box satisfies the
lenient rule trivially and is worthless to a \judge{} that only ever sees the
crop. \figref{fig:proposers} shows one frame where this happens, with the
per-instance IoU printed on the image so the verdict can be checked without
the caption: the photometric stage scores 0.042 on a laptop left on a seat and
0.016 on a bag on the floor --- whose only photometric cover \emph{is} the
98.9\,\% box --- against 0.555 and 0.408 for the feature stage. The figure also
shows what the feature \proposer{} costs: 21 proposals against 8, which is why
\tabref{tab:main} records no saving in judge calls for \armdino{} and why the
veto arm exists at all.

\begin{figure}[t]
  \centering
  \includegraphics[width=\linewidth]{fig2_proposers.jpg}
  \caption{The same inspection frame under both proposal families. Ground truth
  in green, proposals in red, and the box covering 98.9\,\% of the frame
  highlighted in orange. Each ground-truth box carries the best IoU achieved
  and the verdict of the strict rule ($\mathrm{IoU} \ge 0.3$). Both instances
  fail under the photometric stage and pass under the feature stage.}
  \label{fig:proposers}
\end{figure}

\subsection{Twelve judges over a frozen set of boxes}
\label{sec:results:judge}

\tabref{tab:judge} inverts the experiment: the boxes are frozen at the 654
crops produced by \armdinodinomaly{}, and four models are crossed with three
prompts. The calibration anchor is that 73 instances among 654 regions imply
that a calibrated judge answers YES about 11\,\% of the time.

\begin{table}[t]
\centering
\caption{Judge sweep on a frozen set of 654 crops: four vision--language
models $\times$ three prompts, 2\,h\,07 of GPU. Nothing improves on the shipped
configuration at usable specificity.}
\label{tab:judge}
\small
\begin{tabular}{@{}llrrrrr@{}}
\toprule
\textbf{Model} & \textbf{Prompt} & \textbf{YES \%} & \textbf{Instances} &
\textbf{Frame $F_1$} & \textbf{Frame Sp} & \textbf{Reg.\ prec.} \\
\midrule
GLM-V 9B (Flash) & conservative & 10.2 & \textbf{55/73} & \textbf{0.984} & \textbf{1.000} & 0.896 \\
GLM-V 9B (Flash) & balanced     & 10.4 & 55/73 & 0.968 & 0.973 & 0.882 \\
GLM-V 9B (Flash) & lenient      & 13.6 & 56/73 & 0.909 & 0.865 & 0.719 \\
\midrule
InternVL3.5-8B   & balanced     &  8.7 & 48/73 & 0.984 & 1.000 & 0.912 \\
InternVL3.5-8B   & lenient      &  7.8 & 43/73 & 0.931 & 1.000 & 0.922 \\
InternVL3.5-8B   & conservative &  6.9 & 39/73 & 0.893 & 1.000 & 0.933 \\
\midrule
Qwen3-VL-8B      & lenient      & 28.9 & 63/73 & 0.690 & 0.297 & 0.376 \\
Qwen3-VL-8B      & balanced     & 40.5 & 63/73 & 0.659 & 0.189 & 0.268 \\
Qwen3-VL-8B      & conservative & 39.6 & 63/73 & 0.645 & 0.135 & 0.274 \\
\midrule
Qwen2.5-VL-7B    & balanced     & 54.3 & 64/73 & 0.652 & 0.162 & 0.200 \\
Qwen2.5-VL-7B    & conservative & 48.3 & 63/73 & 0.645 & 0.135 & 0.222 \\
Qwen2.5-VL-7B    & lenient      & 88.5 & 63/73 & 0.632 & 0.081 & 0.123 \\
\bottomrule
\end{tabular}
\end{table}

No arm improves on the shipped configuration at usable specificity. That is a
result and not a null result: it is the control that keeps
\secref{sec:results:main} honest, since the first objection to ``the
\proposer{} decides'' is that we chose a weak \judge{}.

Three readings follow. \emph{The 63/73 of the Qwen arms is not detection but
non-rejection}: at an 88.5\,\% YES rate a judge mechanically recovers almost
everything handed to it, and frame specificity collapses to 0.081. This is the
concrete demonstration of why every recall in this literature must be read next
to the YES rate. \emph{The prompt is a real lever only where there is slack}: on
InternVL3.5, which under-fires at 6.9\,\%, the balanced prompt moves 39/73 to
48/73 at specificity 1.000 unchanged; on the GLM-V judge, already calibrated at
10.2\,\%, the same prompt gains zero instances and costs specificity.
\emph{Therefore prompting is calibration rather than capability} --- at this
scale, for this set of models, and without adaptation. All three qualifiers are
necessary: VERA~\cite{Ye_2025} treats prompts as learnable parameters and gains
from it, and Borodin et al.~\cite{Borodin_2025} report that prompt sensitivity
drops sharply after parameter-efficient fine-tuning. Our result and theirs are
two halves of one finding, and the joint statement is that the next lever for
this class of system is adaptation, not further prompt engineering.

\subsection{Are the learned baselines enough?}
\label{sec:results:baselines}

\tabref{tab:baselines} runs anomalib~2.6 reference implementations on the same
cameras, each model fitted on that camera's nominal frames and tested on that
camera's cases. Scores are each method's native image score; I-AUROC is
rank-based, so the comparison is valid despite differing score definitions.

\begin{table}[t]
\centering
\caption{Image-level AUROC per camera, own-camera models. Only the
\texttt{tram\_1762} column carries enough pairs to support a claim
(\secref{sec:setup:sizes}).}
\label{tab:baselines}
\begin{tabular}{@{}lrrrrr@{}}
\toprule
\textbf{Method} & \textbf{3333-c52} & \textbf{3333-c53} & \textbf{3333-c54} &
\textbf{3333-c55} & \textbf{tram\_1762} \\
 & \emph{2 pairs} & \emph{8 pairs} & \emph{8 pairs} & \emph{8 pairs} & \emph{176 pairs} \\
\midrule
PaDiM~\cite{Defard_2021}      & 1.000 & 0.875 & 1.000 & 1.000 & \textbf{0.915} \\
PatchCore~\cite{Roth_2022}    & 1.000 & 0.625 & 1.000 & 1.000 & \textbf{0.841} \\
Dinomaly~\cite{Guo_2025} \emph{(anomalib)}    & 1.000 & 1.000 & 1.000 & 1.000 & \textbf{0.886} \\
AnomalyDINO~\cite{Damm_2025} \emph{(anomalib)} & 1.000 & 0.750 & 1.000 & 1.000 & \textbf{0.835} \\
\midrule
Dinomaly \emph{(ours)}        & 1.000 & 0.750 & 1.000 & 1.000 & \textbf{0.920} \\
Dinomaly2 \emph{(ours)}       & 1.000 & 1.000 & 1.000 & 1.000 & \textbf{0.801} \\
\bottomrule
\end{tabular}
\end{table}

Two things fall out, and the second matters more than the table. Our
re-implementation is sound: it matches the anomalib reference on four of five
cameras, closes the fifth with the Dinomaly2 variant, and on the only column
with real statistics sits slightly above the reference (0.920 against 0.886).
A re-implementation cannot be reported without such a check, and it passes.

More importantly, \emph{frame-level anomaly detection on this data is not the
hard part}. Every standard method, untuned, lands between 0.835 and 0.920 on
the 176-pair column and saturates at 1.000 on three of the four remaining
views. ``Does this frame contain an anomaly'' is close to solved by a 2021
memory-bank baseline that takes nine seconds to fit. This is why the claims of
this paper are made at region and verdict level: \emph{which} region, and
\emph{what} it is --- a language verdict a human can read, which no row of
\tabref{tab:baselines} produces --- and at what operating point that holds.

\subsection{The proposal stage on public data}
\label{sec:results:cdnet}

\tabref{tab:cdnet} evaluates the photometric proposal stage on seven CDnet 2014
categories, without the \judge{} in the loop, so that a pipeline failure can be
attributed: the region was never proposed, or it was proposed and misjudged.

\begin{table}[t]
\centering
\caption{CDnet 2014, photometric proposal stage, pooled over all evaluated
frames. \emph{Mask} is the per-pixel change mask; \emph{boxes} is the
rasterised proposals, which is what reaches the judge. Read the two families of
column separately: the stage is a recall-first proposer by design.}
\label{tab:cdnet}
\small
\begin{tabular}{@{}lrrrrrrrr@{}}
\toprule
\multirow{2}{*}{\textbf{Category}} & \multirow{2}{*}{\textbf{Seq.}} &
\multirow{2}{*}{\textbf{Frames}} &
\multicolumn{4}{c}{\textbf{Mask}} & \multicolumn{2}{c}{\textbf{Boxes}} \\
\cmidrule(lr){4-7}\cmidrule(l){8-9}
 & & & $F_1$ & Re & Pr & FPR & Re & Pr \\
\midrule
\texttt{badWeather}    & 4 &  899 & \textbf{0.675} & 0.557 & 0.856 & 0.0014 & 0.918 & 0.095 \\
\texttt{cameraJitter}  & 4 &  159 & 0.566 & 0.647 & 0.504 & 0.0266 & 0.924 & 0.105 \\
\texttt{shadow}        & 6 &  710 & 0.496 & 0.618 & 0.415 & 0.0463 & 0.942 & 0.178 \\
\texttt{nightVideos}   & 6 &  619 & 0.434 & 0.466 & 0.406 & 0.0100 & 0.886 & 0.129 \\
\texttt{intermittentObjectMotion} & 6 & 1216 & 0.406 & 0.487 & 0.349 & 0.0350 & 0.812 & 0.139 \\
\texttt{dynamicBackground} & 6 & 667 & 0.234 & 0.808 & 0.137 & 0.0714 & 0.998 & 0.028 \\
\texttt{lowFramerate}  & 4 &  264 & \textbf{0.207} & 0.278 & 0.164 & 0.0277 & 0.767 & 0.052 \\
\bottomrule
\end{tabular}
\end{table}

\paragraph{The category that is our task.}
CDnet's \texttt{intermittentObjectMotion} is defined by its authors as videos
containing background objects moving away, abandoned objects, and objects
stopping for a while before moving away --- our task, almost word for word.
Per sequence, the mask $F_1$ ranges from 0.978 on \texttt{streetLight} --- a
fixed camera, stable light, a car that stops, a score that would be competitive
on the CDnet leaderboard --- down to 0.143 on \texttt{winterDriveway} (snow,
changing daylight) and 0.138 on \texttt{tramstop} (outdoor light drift). The
stage is excellent exactly while its reference stays valid and collapses when
it does not, and CDnet finds that failure mode in one afternoon.

\paragraph{Read recall and precision separately.}
Across all seven categories the boxes recall 0.767--0.998 while their precision
stays between 0.028 and 0.178. This is by design, not by accident: our proposal
stage is a recall-first proposer whose false positives are meant to be removed
downstream by a \judge{} that CDnet does not run. Reporting only $F_1$ would
say the stage is weak; reporting both columns says what is true --- it rarely
misses, and it over-proposes on purpose. It should also be noted that the boxes
row is scored against pixel-accurate ground truth by a method that emits
rectangles, which structurally caps its precision.

\paragraph{Which of our assumptions are load-bearing.}
Read top to bottom, \tabref{tab:cdnet} ranks them. \texttt{badWeather} (0.675)
and \texttt{cameraJitter} (0.566) are the good news and are the two categories
that matter most for a tram: snow and rain against a median reference cost
almost nothing (FPR 0.0014), and camera shake --- a tram vibrates continuously
--- is survivable. \texttt{dynamicBackground} (0.234) and
\texttt{lowFramerate} (0.207) are failures, and both lie outside our
deployment: moving water and foliage do not exist inside a vehicle, and we
control the sampling rate. \texttt{shadow} (0.496) and \texttt{nightVideos}
(0.434) are the honest middle, and these two \emph{are} ours: a tram runs at
night and through moving shade. That is the number to quote when asked what the
proposal stage does under illumination change, and it is mediocre.

\paragraph{Two claims of \secref{sec:results:main} reproduce on public data.}
The gated variant \armphotodino{} scores 0.252 mean $F_1$ against the ungated
diff's 0.254, at slightly lower recall (0.898 against 0.921) --- independent
confirmation, on data we did not build, that a veto is a cost lever and never a
quality one. Conversely, the feature \proposer{} does \emph{not} transfer to
CDnet: it reaches 0.283 mean recall and fires not at all on
\texttt{tramstop}. Two explanations were tested and ruled out. Sweeping its
$z$ threshold from 4.0 to 1.5 and its floor from 0.08 to 0.02 moves recall from
0.616 to 1.000 while precision stays pinned between 0.011 and 0.022 and the
false-positive rate climbs to 0.98 --- there is no threshold at which it is
useful; it goes from silent to covering the frame. Rebuilding the reference
from a single real frame rather than a 50-frame median leaves it unchanged
($F_1$ 0.029 against 0.028) while making the photometric stage distinctly
worse, so the composite reference is not the cause. What remains is the
assumption itself: the feature stage compares each patch to the same grid
position of a reference from that same fixed camera, and quantises its boxes to
the patch grid. That is a good trade on a tram interior, where an anomaly spans
a large fraction of a patch, and a bad one on $320\times240$ outdoor footage
scored against pixel-accurate ground truth. \emph{Our best proposal stage is
best conditionally}, and the condition is the deployment geometry.

\subsection{The cost of per-camera models}
\label{sec:results:transfer}

The strongest arm of \tabref{tab:main} needs one trained checkpoint per camera.
Whether that is a genuine cost or an implementation detail depends on whether a
foreign checkpoint actually fails, which we measured across eight cameras.

It fails, and how it fails is the finding. A foreign checkpoint scores a clean
frame $2.13\times$ higher than the correct checkpoint does (median 0.163
against 0.077), and its anomalous-minus-clean margin collapses from
$+0.017 \ldots +0.061$ to a median of $+0.0009$ --- a factor of 41. The
distinction matters more than any AUROC: a model that were merely \emph{offset}
could be rescued by a per-camera threshold, and this one is not offset. On an
unfamiliar camera, ``this view is unfamiliar'' dominates the reconstruction
error and drowns ``this patch is anomalous''. The cause is the viewpoint and
not the illumination: the four \texttt{3333} cameras were filmed in the same
two sessions minutes apart, under identical light, and transfer \emph{within}
that family still fails. One nuance is worth preserving: on the only column
with 176 pairs, foreign checkpoints still reach AUROC 0.665--0.909, so transfer
preserves a weak ordering while destroying the operating point --- exactly the
failure a deployed system cares about and exactly the one an AUROC-only report
would hide.

We then tested the authors' own answer to this objection. Dinomaly2
\cite{Guo_2025_dinomaly2} introduces context-aware recentering (CAR), which
subtracts a global summary of the image from every reconstruction target ---
precisely the global, view-dependent term our measurement implicates. We ported
its four elements and trained eight further checkpoints on identical data,
crops, schedule and holdout, then ablated CAR alone.
\tabref{tab:dinomaly2} attributes the effect cleanly. Because CAR normalises
the reconstruction targets, absolute score levels are not comparable across
variants and only the ratios are; every entry in the first two columns is a
ratio for that reason.

\begin{table}[t]
\centering
\caption{Cross-camera ablation. Context-aware recentering is the only element
that touches the margin --- and it is also what costs the home camera.}
\label{tab:dinomaly2}
\begin{tabular}{@{}lrrrr@{}}
\toprule
\textbf{Variant} & \textbf{Clean inflation} & \textbf{Margin collapse} &
\textbf{Transfer} & \textbf{tram\_1762} \\
 & \emph{foreign\,$\div$\,own} & \emph{own\,$\div$\,foreign} &
 \emph{median, pooled} & \emph{self, 176 pairs} \\
\midrule
Dinomaly (v1)          & $2.13\times$ & $41\times$ & 0.367 & \textbf{0.920} \\
Dinomaly2, CAR off     & $1.73\times$ & $43\times$ & 0.393 & \textbf{0.903} \\
Dinomaly2, CAR on      & $\mathbf{1.52\times}$ & $\mathbf{10\times}$ & 0.449 & \textbf{0.801} \\
\bottomrule
\end{tabular}
\end{table}

Without CAR the margin collapse is $43\times$, statistically indistinguishable
from v1's $41\times$; with it, $10\times$. The loose loss, the bottleneck shape
and the reversed decoder pairing do nothing at all for cross-camera behaviour
--- what they fix is the diagonal, where \texttt{3333-cam53} goes from 0.750 to
1.000 with CAR disabled. The mechanism proposed above is therefore confirmed by
removing exactly the term it implicates and watching the effect disappear. CAR
is also what costs the home camera: 0.903 without it, essentially v1's 0.920,
against 0.801 with it. Recentering subtracts a global term carrying nuisance
\emph{and} signal, and on the camera where we have the most data that trade is
negative.

The honest summary is a trade-off rather than an improvement: context-aware
recentering buys a fourfold better cross-camera margin and pays about 0.10
image-level AUROC at home, and even the improved margin leaves transfer at
chance. Per-camera training is therefore a measured requirement of this design,
not a cost we chose --- and to our knowledge it is not reported by the works
citing either Dinomaly or AnomalyDINO, whose benchmarks contain no notion of a
viewpoint that a model could fail to transfer across.

\subsection{Qualitative check on abandoned objects}
\label{sec:results:aboda}

ABODA is the one public dataset whose anomaly class is ours, and it ships no
annotations, so \figref{fig:aboda} is a figure and not a number. It shows the
proposal stage on an outdoor night sequence it was never tuned for: while
people are present the stage proposes them, and once they leave, the bag they
left behind is proposed alone. The behaviour is the one the design intends in a
domain --- outdoor, night, crowded --- that shares nothing with a tram
interior except the observation geometry.

\begin{figure}[t]
  \centering
  \includegraphics[width=0.7\linewidth]{fig4_aboda.jpg}
  \caption{Proposal stage on ABODA, qualitative. The reference is the median of
  the opening frames and the person veto is disabled, as on CDnet. Once the
  people have left, the abandoned bag is the only proposal.}
  \label{fig:aboda}
\end{figure}
